This paper has demonstrated that the predictive power of the yield curve for US recessions exhibits fundamental regime dependence, varying systematically with the nature and origin of economic contractions. Through comprehensive analysis spanning four decades and employing both traditional econometric methods and modern machine learning techniques, we have established three central findings that extend the foundational work of \citet{estrellaYieldCurvePredictor1996} into the era of unconventional monetary policy.

First, yield curve factors extracted via Principal Component Analysis, particularly the slope factor, retain strong predictive power during endogenous financial crises such as the 2007-2009 Global Financial Crisis, achieving balanced accuracy above 0.83 when paired with regularized linear models. The inversion of the yield curve continues to serve as a reliable early warning signal when recessions emerge from monetary tightening cycles and financial sector stress, consistent with the traditional interpretation of yield spreads as reflecting market expectations of future policy easing. Second, this predictive relationship breaks down almost entirely during exogenous shocks exemplified by the COVID-19 pandemic, where real-sector macroeconomic indicators such as unemployment and industrial production provide superior forecasting performance with balanced accuracy reaching 0.86. Third, factor loading analysis reveals substantial structural instability in yield curve dynamics during periods of unconventional monetary policy, with pronounced drift in the importance of level, slope, and curvature factors around quantitative easing programs and zero lower bound episodes, confirming the concerns raised by \citet{rudebusch2006macroeconomic} regarding term premium distortions.

These results carry important implications for both academic research and practical policy analysis. The failure of yield curve models during COVID-19 underscores the necessity of adaptive forecasting frameworks that can distinguish between monetary-policy-driven slowdowns and sudden external disruptions. Our findings suggest that robust early warning systems should implement regime-switching ensemble architectures that dynamically adjust the weight assigned to financial versus real-sector indicators based on leading indicators of crisis type. Furthermore, the documented distortions in term premia during unconventional policy periods support the use of adjusted slope measures that isolate expectations components from non-expectational yield movements, as advocated by \citet{sharpeNearTermForwardYield2018}.

Future research should extend this work in several directions. The development of formal regime classification algorithms capable of real-time crisis-type identification would enable more systematic application of our findings. Integration of higher-frequency data sources, including weekly jobless claims, credit card spending, and mobility indicators, could improve forecast performance during rapid-onset recessions where monthly macroeconomic data lag. Additionally, incorporating formal term structure decompositions would allow for cleaner isolation of risk-neutral expectations from term and liquidity premia. Finally, expanding the analysis to international contexts would test whether the regime-dependent patterns we document for the United States generalize across different institutional and monetary policy frameworks.